\section{Introduction}
\label{sec:intro}

% TODO: full intro. Sketch:
%
% - BCI as a research field has shifted from clinical EEG (PSG, multi-thousand
%   dollar systems) to consumer-grade hardware (Muse, OpenBCI, etc.). This
%   democratizes the field but introduces a new engineering challenge: the
%   acquired signal is noisier, the channel count is smaller, and the
%   scientific claims that hold for clinical PSG do not always transfer.
%
% - Consumer-grade EEG has been used in three application areas: (a) BCI for
%   accessibility and communication (the work of the present platform); (b)
%   meditation and wellness; (c) sleep tracking. This work is in (a) and
%   aims to extend to (c) with the same hardware.
%
% - Sleep-stage estimation from consumer-grade hardware is a known challenge.
%   Recent work (cite) has shown that 4-channel frontal EEG can estimate
%   sleep stage with moderate accuracy. The present work builds on this
%   literature.
%
% - AI-assisted dream incubation is a recent research direction. The Dormio
%   system \citep{maes2020dormio} demonstrated that hypnagogic-state audio
%   cues can incubate dream content. The present work extends this to a
%   fully automated, on-device pipeline.
%
% - The contribution of this paper is the platform, not the empirical
%   question. We describe the architecture (§\ref{sec:methods}), the
%   validation (§\ref{sec:results}), and the planned study (§\ref{sec:methods}-D8).
%   We do not claim that the platform improves creative problem solving.
%   That is a pre-registered empirical question to be tested in future work.

The NeuralCompose platform is a privacy-first, fully on-device macOS
research artifact for EEG-driven communication and sleep-cycle research.
This paper describes the platform, the validation that the
hardware-and-software pipeline is operational, and the pre-registered
research program that will test whether AI-assisted dream incubation
improves creative problem solving.

The platform is implemented in Swift 6 with strict concurrency. EEG
acquisition uses BrainFlow \citep{brainflow2022} over native Bluetooth Low
Energy on the macOS host's Core Bluetooth stack. Sleep-stage classification
runs as a small convolutional neural network on the Apple Neural Engine
(ANE) via Core ML. The local large language model (LLM) uses the MLX-Swift
runtime. No network calls occur at any point during a session.

The platform is open source at \url{https://github.com/aurascoper/NeuralCompose}.
The repository contains the full type-level design
(\texttt{SLEEP\_CYCLE\_DESIGN.md}, $\sim 1{,}500$ lines), the validation
toolkit (\texttt{Sources/BCIEEG/EEGScalpPlotterView.swift},
\texttt{Scripts/validate-muse-physiology.py}), and a draft manuscript
(\texttt{paper/}). The platform is useful independently of the empirical
question.

The remainder of the paper is organized as follows. Section
\ref{sec:methods} describes the architecture, the validation protocol, and
the planned experimental study. Section \ref{sec:results} reports the
current validation results. Section \ref{sec:discussion} discusses the
implications, the limitations, and the open research questions.
